\chapter{Chapter 8 --- Discussion}\label{chapter-8-discussion}

Five phases were designed to answer four questions in sequence. In the
event, two returned answers the design did not anticipate, one returned
a negative, and the last overturned the explanation for its own positive
result. This chapter reads them as a single argument, states what
follows for researchers and for practitioners, and gives the full
treatment of validity threats.

\section{8.1 Synthesis across the research
questions}\label{synthesis-across-the-research-questions}

\subsection{8.1.1 The four answers}\label{the-four-answers}

\textbf{RQ1 --- How much do SZZ variants disagree, and in which
direction?} Substantially, and structurally. No variant exceeds
\textbf{27.2\% precision} against the reference. Variants agree with one
another at \ensuremath{\kappa} up to \textbf{0.933} while agreeing with the reference at \ensuremath{\kappa}
\textbf{0.158--0.202} --- four to six times more strongly with each
other than with the thing they are estimating. Bias is bifurcated: B-SZZ
over-flags (\ensuremath{\rho}\ensuremath{_0} = 0.263), L-SZZ and R-SZZ under-flag (\ensuremath{\rho}\ensuremath{_1} = 0.733, 0.700).
\textbf{Two decades of refinement have relocated error rather than
reduced it}, moving it from the false-positive column to the
false-negative column.

\textbf{RQ2 --- How does label source shift measured performance across
regimes?} It shifts it significantly --- 0.034 to 0.064 MCC under
realistic streaming, against a deployable baseline near 0.097.
\textbf{But the evaluation protocol shifts it far more.} Random
\emph{k}-fold inflates a random forest by \textbf{+0.137 MCC}; scoring
against the training labels inflates it by a further \textbf{+0.230}.
The two protocol effects together are roughly nine times the
label-source effect.

\textbf{RQ3 --- Through what mechanism does noise degrade an online
learner?} Through amplification rather than starvation. Repairing
B-SZZ's false positives recovers \textbf{+0.040 MCC}; restoring its
false negatives is \emph{equivalent to no repair} (TOST \emph{p} =
0.0008) under realistic, immediate and at-window delivery alike. The
mechanism is measurable: ORB's oversampling rate is a near-perfect
inverse of delivered defect-label supply (\textbf{\ensuremath{\rho} = \ensuremath{-}1.000}). The
boost compensates for label \emph{scarcity} and cannot distinguish a
scarce-but-correct stream from a scarce-and-wrong one.

\textbf{RQ4 --- Can a learner recover that loss without the reference?}
Not by identifying wrong labels. Phase 4's filter is adoptable --- the
non-degradation gate passes --- but every confirmatory test points the
predicted way and \textbf{none survives correction}; the oracle-assisted
repair recovers +0.040 against the approximation's +0.017. Phase 5 then
registered the damping alternative on new data and confirmed it beats
both baseline and filter, \textbf{but a control showed the entire
benefit is the absence of oversampling-rate boosting rather than any
defence}: plain online bagging matches every noise-aware variant and
beats the boosted learner in \textbf{18 of 18 projects} under
false-positive-heavy noise. So the answer is \emph{no} to identifying
wrong labels and \emph{yes} to something cheaper --- stop amplifying
them.

\subsection{8.1.2 The noise pipeline, read end to
end}\label{the-noise-pipeline-read-end-to-end}

The four results compose into one account of how a number becomes
untrustworthy.

\textbf{Stage 1 --- The labels are wrong, asymmetrically.} Tangled
commits cause SZZ to blame lines that fix nothing. Filters remove some
of that over-collection and, for the aggressive variants, more than half
of what they remove is correct.

\textbf{Stage 2 --- The evaluation hides it.} Because the labels are
heuristic, models are scored against the same heuristic that trained
them. A high-capacity model that fits the heuristic's error structure is
rewarded for doing so. \textbf{This is only possible because the labels
are wrong}: with ground truth there would be nothing to score against
but truth, and the +0.230 gap could not exist. The two problems are not
independent --- the labelling failure creates the evaluation failure.

\textbf{Stage 3 --- Deployment removes the protections.} Under streaming
with verification latency, the learner's imbalance machinery amplifies
whatever defect labels arrive. Most of the surviving labels are false
positives, and they are amplified hardest exactly when defect labels are
scarcest --- which, given that 53\% arrive after the decision window, is
always.

\textbf{Stage 4 --- Knowing this is not sufficient to fix it, but it is
sufficient to avoid it.} Identifying \emph{which kind} of error hurts
does not confer the ability to identify \emph{which instances} are that
error. A confidence signal cannot separate ``this label is wrong'' from
``I have not learned this pattern yet'', and under an 8.5\% positive
rate the second case is common. What \emph{is} available is the
amplifier itself: the learner cannot find the wrong labels, but it can
decline to magnify them. \textbf{Imbalance correction and label noise
interact, and the standard remedy for the first makes the second worse}
--- which is the one actionable result the five phases produce.

\textbf{The unifying quantity is deliverable precision.} What an online
learner needs is not correct labels in the abstract but correct labels
that \emph{arrive} --- and among those that arrive, false positives are
the error it cannot defend against, because its own imbalance mechanism
amplifies them. Every result in this thesis is a corollary of that
sentence, including the retirement of the starvation hypothesis it began
with.

\subsection{8.1.3 Statistical significance and practical
significance}\label{statistical-significance-and-practical-significance}

The measurement critique is statistically overwhelming: rank-biserial
+1.000 on the leakage effect (21 of 21 projects), +0.991 on the
self-scoring gap, intervals excluding zero by wide margins, survival
under global correction.

\textbf{Practical significance is more modest and should be stated as
such.} Even with perfect labels and an honest protocol, deployable MCC
on this corpus is approximately \textbf{0.097}. Repairing every false
positive raises it to 0.096 from 0.058 --- a 66\% relative gain on a
small absolute base. \textbf{This thesis does not claim that better
labels would make JIT-SDP good.} It claims they would make it measurably
less self-sabotaging, and that the field's reported numbers overstate
what it currently achieves by roughly a factor of four.

The low ceiling is itself a finding. A field reporting 0.41 and
delivering 0.10 has a calibration problem independent of whether 0.10 is
useful.

\section{8.2 Implications}\label{implications}

\subsection{8.2.1 For researchers building
datasets}\label{for-researchers-building-datasets}

\textbf{Stop treating inter-tool agreement as validation.} Variants
agreeing at \ensuremath{\kappa} 0.93 while agreeing with the reference at \ensuremath{\kappa} 0.16 is the
clearest result in Chapter 4. Agreement between procedures that share a
blame step and a set of assumptions measures shared failure modes.

\textbf{Report the denominator.} A variant scored on the candidates it
emits will appear more precise than one scored on the whole corpus.
Chapter 4's comparison is only interpretable because every variant was
scored over an identical 27,319-commit population.

\textbf{Report both flip rates, not a single noise rate.} \ensuremath{\rho}\ensuremath{_0} and \ensuremath{\rho}\ensuremath{_1}
differ by a factor of eleven across the variants studied here. Any
downstream noise model calibrated on a single scalar will be wrong for
five of the six.

\textbf{Distinguish rate from volume.} At an 8.5\% positive rate, a \ensuremath{\rho}\ensuremath{_0}
of 0.263 produces 6,565 wrong labels while a \ensuremath{\rho}\ensuremath{_1} of 0.733 produces 1,710.
Chapter 6's central finding turns entirely on this distinction, and the
error-matched control that established it is cheap to run.

\subsection{8.2.2 For researchers reporting
results}\label{for-researchers-reporting-results}

\textbf{Random \emph{k}-fold on temporally ordered data is not
defensible}, and the inflation it produces is \emph{larger for models
with more capacity to exploit it}. Comparisons between model families
measured this way are confounded, and the confound flatters exactly the
models the field most wants to promote.

\textbf{Score against something other than the labels you trained on},
or state plainly that the number measures fidelity to a heuristic. Where
no independent reference exists, the honest description of a self-scored
result is ``reproduces SZZ at MCC \emph{x}'', not ``detects defects at
MCC \emph{x}''.

\textbf{Name the prequential estimator.} A terminal fading value and a
trajectory mean disagree on the sign of one comparison in this thesis
and would have supported a claim that had to be withdrawn. Neither is
canonical for MCC, which is not a decomposable loss.

\textbf{Report multiplicity correction globally when families were not
pre-declared.} The distinction costs one column and removes an entire
class of objection.

\subsection{8.2.3 For practitioners}\label{for-practitioners}

\textbf{Expect roughly 0.10 MCC, not 0.40.} If a JIT-SDP tool is
evaluated on your own repository under an honest protocol and returns
something near 0.10, it is performing as well as the state of the
practice on this corpus --- not failing.

\textbf{Label quality is worth more than model sophistication.}
Reference labels beat every SZZ variant under streaming; a
single-feature logistic regression beats a 100-tree forest under a
chronological split. Investment in knowing which commits actually
introduced defects dominates investment in the learner.

\textbf{If you must use SZZ, prefer the permissive variant and accept
the false positives} --- B-SZZ carries the smallest penalty of the six
under streaming, and the aggressive filters discard more correct answers
than incorrect ones. This runs against the direction of variant
development and is stated because the data support it.

\textbf{Switch off prediction-bias boosting when your labels come from
SZZ.} This is the clearest practical result in the thesis: plain online
bagging beat the boosted variant in 18 of 18 projects under
false-positive-heavy noise, matched every noise-aware method tested, and
costs nothing to implement. The boost is designed for clean imbalanced
streams, and SZZ-labelled streams are not clean.

\textbf{Do not deploy a confidence-based label filter or damper on the
strength of this work.} Chapter 7 tested both. Each passed the
non-degradation gate; neither improved on simply removing the boost.

\section{8.3 Threats to validity}\label{threats-to-validity}

\subsection{8.3.1 Construct validity}\label{construct-validity}

\textbf{The reference is not ground truth, and every comparison inherits
that.} It is \texttt{git\ blame} seeded with human-verified fix lines.
Comparisons against it isolate the cost of tangled commits, not total
labelling error. \textbf{All measured noise is therefore a lower bound}:
blame error is present on both sides and cancels.

\textbf{The reference condition's label timing is SZZ-derived.} This is
the most serious threat in the study. The published labels carry no
native fix-to-inducing linkage, so arrival times were reconstructed from
the union of SZZ mappings. The reference is blame-derived in
construction \emph{and} SZZ-derived in timing. A coverage sensitivity
analysis (67.8\% \ensuremath{\rightarrow} 100\%) bounds part of the exposure; the timing itself
is not bounded, and every streaming result in Chapters 5--7 inherits it.

\textbf{The corpus is SZZ-shaped.} Changes adding no new lines were
excluded by the dataset authors, so defects introduced purely by
deletion cannot appear. Separately, because the reference is blame
output, every reference positive is blame-reachable by construction, so
blame-unreachability is not an available explanation for any measured
false negative. This limit was declared before Phase 3 ran.

\subsection{8.3.2 Internal validity}\label{internal-validity}

\textbf{Phase 3's repair conditions are not error-matched, and the
matched control changed the claim.} Repairing ``all of each'' compares
6,565 corrections against 837. At equal mass the two error types are
indistinguishable. The supported claim is about error \emph{volume}; the
per-label claim is explicitly disclaimed.

\textbf{Phase 3's dose parameterisation is not class-balanced.} Dose is
a fraction of all commits, so the FN-heavy profile strips far more of
the minority class at equal dose, and at the highest dose leaves a
stream whose defect labels are \emph{all false}. Slopes are fitted only
on cells retaining genuine defect labels.

\textbf{Two published conclusions were wrong and are corrected rather
than removed.} An earlier version reported that latency does not
compress sensitivity to label quality; that comparison ran between two
uniformly delayed arms and could not test the question, and against a
true no-latency control it reverses. Separately, Phase 4's exploratory
analysis was read as evidence that confidence damping works; Phase 5's
control shows the benefit is attributable to removing the boost, with
damping adding nothing over plain online bagging in any noise profile.

\textbf{Phase 5's registration contained a mis-specified test.} The
non-degradation gate was registered as an equivalence (TOST) test, which
returns ``not equivalent'' because damping is \emph{better} than the
baseline by more than the margin. Non-inferiority is the correct test
for that gate. The direction is unambiguous either way; the error is
reported because it was registered.

\textbf{The latency decomposition is unresolved, not null.} A full 2\ensuremath{\times}2
was built after the original decomposition was found to be unidentified.
Every contrast, including the interaction, has an interval spanning
zero.

\textbf{A label-vintage defect was found and fixed mid-study.} Models
had trained on labels one commit older than the noise rates reported,
affecting 1.0--6.5\% of commits per variant. The remediation ships as an
executable gate that runs before any compute.

\textbf{Phase 4's held-out projects were not randomly assigned.} They
were chosen after Phases 1--3 were analysed. Nothing in Phase 4's
reported numbers depends on them, but the assignment was purposive.

\subsection{8.3.3 External validity}\label{external-validity}

\textbf{One benchmark family.} Twenty-one Apache Java repositories of a
single dataset lineage. Java-specific idioms, Apache-specific review
culture and the dataset's own inclusion criteria are all uncontrolled.

\textbf{Project heterogeneity is real but does not drive the results.}
Projects span 544--4,026 commits and 19--335 defects and are weighted
equally. Excluding those below a 25-defect floor moves no headline
estimate by more than 0.002 MCC and all four Tier 1/Tier 2 claims still
survive correction.

\textbf{The mechanism is architecture-specific.} The amplification
account applies to learners whose oversampling responds to observed
class rate. A learner without that machinery would have a different
failure mode, and Chapter 6's conclusions should not be transferred to
one.

\subsection{8.3.4 Conclusion validity}\label{conclusion-validity}

\textbf{Effects below roughly 0.04 MCC are not resolvable at \emph{n} =
21.} This is load-bearing for the latency decomposition and for Phase
4's H1 (+0.017), and both are reported as unresolved rather than as
null.

\textbf{The exploratory comparison count is large}, which is why global
correction is reported alongside within-family correction and why
headline claims are those surviving the global column.

\textbf{Phase 4's negative could be a power failure rather than an
absence.} Four registered tests agreeing in direction, with H1's
interval barely excluding zero, is consistent with a real effect too
small to resolve at \emph{n} = 18. The chapter says ``not demonstrated
to work'', not ``shown not to work'', and the distinction is deliberate.

\subsection{8.3.5 Threats that were anticipated versus
discovered}\label{threats-that-were-anticipated-versus-discovered}

Anticipated in Chapter 3 and unchanged: the reference's construct
limits, the single benchmark family, the equal weighting of unequal
projects, the resolution floor.

\textbf{Discovered during the work}, and recorded because the pattern is
itself informative: the error-matching confound, the dose
parameterisation trap, the degenerate quantile rule, the label-vintage
cache defect, the unidentified latency decomposition, and the no-latency
arm that was not one. Every one of these was found by a control that had
been built for a different purpose, or by a reviewer asking what a
caveat was doing in place of an experiment. \textbf{A caveat is not a
control}, and five of the six corrections in this thesis exist because
that distinction was eventually enforced.
